\appendix
\section{Development history and additional results}
\label{app:history}

\paragraph{Two fixes made during development.} In the interest of transparency, we report two problems found while building the pipeline. Both were diagnosed on the development split and fixed before the final run reported in the main text.
\begin{enumerate}
\item \textbf{Mixed-anion minerals.} Our first chemistry rule labeled any RRUFF species containing CO$_3$ (or SO$_4$) as a carbonate (or sulfate), even if it also contained another oxyanion group. As a result, arsenate-bearing species such as beudantite and tyrolite were counted as sulfates or carbonates. Their strongest bands sit at 800--870\,cm$^{-1}$, on top of olivine's doublet, so they ``won'' many Mars olivine spectra. The corrected rule counts a species as a carbonate or sulfate only if it contains no other anion-forming element. The number of sulfate species fell from 191 to 141 and carbonate species from 138 to 111.
\item \textbf{How species probabilities become a group.} We first summed the CNN's species probabilities within each group. Under domain shift, a CNN spreads its probability across many species, so summing favors groups with many species (141 sulfates against 7 olivines). The final rule takes the group of the single most probable allowed species, the same rule library search uses. For the clean-lab CNN, summing gives \sumBalCnnCleanMars{} balanced accuracy, against \balCnnCleanMars{} with the final rule. For the CNN with simulated Mars and the fiber-bump template, the two rules give \sumBalCnnMarsifiedArtifactMars{} and \balCnnMarsifiedArtifactMars{}.
\end{enumerate}
An earlier version of the pipeline that contained both problems was also scored once on the test split. Its numbers are not reported because they reflect these errors. All configurations of the final pipeline that we evaluated are reported in Table~\ref{tab:main}.

\paragraph{Development-split scores.} On the \nDevEval{} development spectra (Mars-candidate setting), balanced accuracy was \devCosine{} for library search, \devCosineWindow{} for library search with the bump skipped, \devCnnClean{} for the clean-lab CNN (the CNN selected for the main text), \devCnnMarsified{} for the simulated-Mars CNN, and \devCnnMarsifiedArtifact{} for the simulated-Mars CNN with the template. The development split has few classes (no feldspar or silica), so these numbers are noisier than the test numbers. Still, both splits rank library search with the bump skipped first and the clean-lab CNN second, and both rank the simulated-Mars CNN last among these five.
